% % University of Minnesota Beamer Template
%\documentclass[handout]{beamer} % Uncomment to get handouts
\documentclass[11pt,xcolor=dvipsnames, aspectratio=169]{beamer}

\usepackage{shyne}
\usepackage{graphicx}
\newcommand{\bc}[1]{\bt{\texttt{#1}}}

% % % Fonts: Feel free to comment/uncomment as needed
%\usepackage{newcent}
\usepackage{palatino}
%\usepackage[T1]{fontenc}
%\usepackage[scaled]{helvet}

\mode<presentation>{%

% Define the ``official'' UMN colors
\xdefinecolor{GopherMaroon}{HTML}{7A0019}
\xdefinecolor{GopherGold}{HTML}{FFCC33}
\xdefinecolor{GopherLightGold}{HTML}{FFDE7A}
\xdefinecolor{GopherDarkMaroon}{HTML}{5B0013}

\usecolortheme[named=GopherMaroon]{structure}
%\usecolortheme[named=GopherGold]{structure}

\usecolortheme{rose} % Try 'lily' if you want to get rid of block backgrounds
\useoutertheme[subsection=false]{smoothbars}
\setbeamertemplate{headline}{} % Uncomment to get standard 'smoothbars' headers back
\setbeamertemplate{footline}{}

\useinnertheme{rounded}
\usefonttheme{structurebold}
\usefonttheme[onlymath]{serif}


\setbeamertemplate{items}[circle]
\setbeamertemplate{sections/subsections in toc}[circle]
\setbeamertemplate{note page}[plain2]

\setbeamercolor{frametitle}{fg=GopherGold,bg=GopherMaroon}
\setbeamercolor{section in head/foot}{fg=GopherMaroon,bg=GopherGold}


% % % % Get rid of the navigation symbols
\setbeamertemplate{navigation symbols}{}

% % % Draw the UMN wordmark in the lower right-hand corner
\logo{\includegraphics[width=1.2in]{../images/maroonWM.png}}
}



\title[Short title]{Version Control Concepts (with Git)}
\author[M. Shyne]{\Large Michael Shyne}
\institute[Short Inst]{Biostatistical Design and Analysis Center\\University of Minnesota}
\date{October 22, 2018}

\begin{document}

\frame{\titlepage}

\begin{frame}{Why use version control?}
There are two main purposes for version control:
\begin{itemize}
\item To maintain an historical record of changes to code, documents or data
\begin{itemize}
\item To be able to easily undo changes that led to errors or other problems
\item To be able to return to a version of particular importance (e.g published paper)
\end{itemize}

\item To manage collaboration between a team of developers
\end{itemize}
\end{frame}

\begin{frame}{Basic concepts of version control}
\begin{itemize}
\item Version control systems (like Git) maintain \bt{repositories}, stores of documents and their changes over time. Repositories can be local (located on your own computer) or remote (located at a sharable site, like GitHub).

\item A version of a document is added to a repository by performing a \bt{commit}. A commit should always be accompanied by a short description of the changes made in this version of the document.

\item A \bt{push} will update a remote repository with changes committed to a local repository.

\item Important milestones can be marked by \bt{tag}.

\item A \bt{branch}, or \bt{fork} is a copy of a repository made a particular point in time. Changes can be committed to a branch without affecting the original (\bt{trunk}).

\item A branch can be reintegrated with the trunk by a \bt{merge} or \bt{pull}.
\end{itemize}
\end{frame}

\begin{frame}{Important git commands}
\begin{itemize}
\item \bc{git init}: Create a local repository
\item \bc{git clone}: Create a local copy of a remote repository
\item \bc{git add}: Stage a change, i.e. mark a file to be committed
\item \bc{git commit}: Add all staged documents to repository
\item \bc{git status}: Shows the status of changed files, could be untracked, modified or staged. 
\item \bc{git push}: Updates remote repository with commits made to local repository (since last push)
\item \bc{git pull}: Updates local repository with commits made to remote repository
\item \bc{git merge}: Combines two branches
\end{itemize}
\end{frame}

\end{document}